\section{Results}\label{sec:results}

\subsection{Main Results: 5-Day Forecast Horizon}

\Cref{tab:main_results} reports the out-of-sample performance of the seven model configurations (six individual models plus the ensemble) for the primary forecast target, 5-day forward realized volatility, over the test period January 2022 through November 2025.

\begin{table}[H]
\centering
\caption{Out-of-sample volatility forecasting performance (5-day horizon, 2022--2025, $N = 980$ trading days). QLIKE and RMSE: lower is better. MZ $R^2$: higher is better. Best value in each column is in \textbf{bold}.}
\label{tab:main_results}
\begin{tabular}{l c c c c}
\toprule
\textbf{Model} & \textbf{QLIKE} & \textbf{RMSE} & \textbf{MAE} & \textbf{MZ $R^2$} \\
\midrule
GARCH(1,1)   & 0.3806 & 0.0796 & 0.0523 & 0.2835 \\
EGARCH(1,1)  & 0.3748 & 0.0769 & 0.0518 & 0.3097 \\
GJR-GARCH    & 0.3447 & \textbf{0.0734} & 0.0489 & \textbf{0.3802} \\
\midrule
HAR-RV       & 0.4198 & 0.0779 & 0.0507 & 0.2895 \\
\midrule
LightGBM     & 0.3632 & 0.0742 & 0.0476 & 0.3754 \\
XGBoost      & 0.3553 & 0.0745 & \textbf{0.0470} & 0.3638 \\
\midrule
Ensemble     & \textbf{0.3431} & 0.0738 & 0.0475 & 0.3515 \\
\bottomrule
\end{tabular}
\end{table}

Several patterns are worth highlighting.

\textbf{GJR-GARCH is the best individual model on the full test sample.} This was not what we expected going in. The asymmetric GARCH specification, which gives negative returns a larger weight in the variance equation, outperforms both tree-based methods and the standard GARCH on QLIKE (0.3447 vs.\ 0.3806), delivers the lowest RMSE (0.0734), and achieves the highest Mincer-Zarnowitz $R^2$ (0.3802). The leverage effect, it turns out, is not just a well-known empirical regularity; it is one of the most valuable pieces of information for forecasting volatility over this test period. The 2022 bear market, with its sharp, asymmetric drawdowns, rewarded models that respond aggressively to negative returns. We caution that this full-sample ranking masks substantial regime dependence, which we examine in the subperiod analysis below and in \Cref{sec:robustness}.

\textbf{Tree-based methods are strong but do not dominate.} LightGBM and XGBoost both outperform standard GARCH and EGARCH, confirming that conditioning on a richer feature set (VIX, lagged RV at multiple horizons, leverage indicators) improves forecasts. XGBoost edges out LightGBM on QLIKE (0.3553 vs.\ 0.3632) and produces the lowest MAE in the table (0.0470, narrowly ahead of GJR-GARCH's 0.0489), though the differences are small. On QLIKE and RMSE, however, neither tree model beats GJR-GARCH. This suggests that the asymmetric volatility response that GJR-GARCH captures parametrically with a single indicator term is the dominant signal, and the 35-feature panel available to the tree models adds only marginal value on top.

\textbf{HAR-RV underperforms, and the regime-conditional analysis in \Cref{sec:robustness} explains why.} The three-parameter linear model, which has been remarkably competitive in prior studies, ranks last on full-sample QLIKE (0.4198). The regime-conditional decomposition in \Cref{tab:subperiod_regime} shows that the underperformance is concentrated in the high-RV quartile, where HAR-RV's QLIKE balloons to 0.7140 because a linear regression on past RV cannot extrapolate to days of extreme realized volatility. In the lower-RV quartiles HAR-RV is competitive, with QLIKE 0.3217 (third best). HAR conditions only on past realized volatility and has no access to VIX or to asymmetric response terms, so when the RV level moves outside the convex hull of the training distribution the linear projection breaks. This is a known weakness of OLS on non-stationary RV; the 5-minute realized-volatility variants of HAR studied in the high-frequency literature partially address it, but require intraday data that we deliberately do not use here.

\textbf{The ensemble wins on QLIKE.} The equal-weight average of LightGBM, HAR-RV, and GARCH(1,1) achieves the lowest QLIKE of any model (0.3431), edging out GJR-GARCH by a narrow margin. This result is consistent with the forecast combination literature \citep{Timmermann2006}: combining structurally diverse forecasters (a parametric time-series model, a linear realized-volatility model, and a nonparametric ML model) smooths idiosyncratic errors. The ensemble does not have the highest MZ $R^2$ (GJR-GARCH does), which means the ensemble is better calibrated on average but slightly less correlated with realized volatility point-by-point. For risk management applications where calibration matters more than correlation, the ensemble is the safer choice.

\subsection{Feature Importance}

\Cref{tab:feature_importance} reports the top ten features by LightGBM split-based importance from a single fit on the full training+validation window (2004--2021), with early stopping on the last 20\% of that window. The model converged at 80 boosting rounds.

\begin{table}[H]
\centering
\caption{Top 10 features by LightGBM split-importance (single fit on 2004--2021 train+validation window, $n_{\text{boost}} = 80$). Share is split count divided by total split count across all 35 features.}
\label{tab:feature_importance}
\begin{tabular}{c l r r}
\toprule
\textbf{Rank} & \textbf{Feature} & \textbf{Splits} & \textbf{Share} \\
\midrule
1  & VIX close (lag 1)                         & 213 & 8.9\% \\
2  & Daily log return (lag 0)                  & 198 & 8.3\% \\
3  & Negative return, 5-day cumulative         & 154 & 6.4\% \\
4  & 5-day RV (lag 22)                         & 148 & 6.2\% \\
5  & HAR monthly component                     & 147 & 6.1\% \\
6  & VIX, 5-day average                        & 145 & 6.0\% \\
7  & Volume / 22-day MA ratio                  & 131 & 5.5\% \\
8  & 5-day RV (lag 5)                          & 113 & 4.7\% \\
9  & Positive return, 5-day cumulative         & 108 & 4.5\% \\
10 & 5-day RV (lag 10)                         & 102 & 4.3\% \\
\bottomrule
\end{tabular}
\end{table}

Three observations on this ranking. First, VIX is the most-split feature. The VIX is the market's forward-looking expectation of 30-day volatility, derived from S\&P 500 option prices, and aggregates information about upcoming FOMC meetings, earnings season, and geopolitical events that no backward-looking feature can capture. That LightGBM identifies VIX as its most valuable input is sensible: the model is effectively learning to combine a market-implied forecast with statistical features from past data. The 5-day VIX average (rank 6) provides a smoother version of the same signal.

Second, the leverage effect is captured through three different features: today's log return (rank 2), the cumulative negative return over the past five days (rank 3), and the cumulative positive return over the past five days (rank 9). Together these three features account for roughly 19\% of total split count, suggesting the tree model spends a substantial fraction of its capacity learning the same asymmetric response that GJR-GARCH encodes structurally with two parameters.

Third, the long-horizon RV signals dominate among the lagged-RV features (lag 22 at rank 4, the HAR monthly component at rank 5, lag 10 at rank 10), while the daily and weekly HAR components do not crack the top ten. The model treats short-horizon RV as already absorbed by the more recent return-based features, and uses the long-horizon channels to anchor the forecast level. Volume relative to its 22-day moving average (rank 7) earns a place in the top ten, consistent with the practitioner intuition that abnormal volume precedes abnormal volatility, though it is not as dominant as the price-based features.

\subsection{Subsample Analysis: A Surprising Reversal}

The full-sample ranking masks substantial regime dependence. To examine this, we split the test period two ways: by calendar year (2022 standalone vs.\ 2023--2025) and by volatility regime (top-quartile RV days vs.\ the remaining three quartiles). \Cref{sec:robustness} reports the full numerical table; the headline pattern is striking enough to deserve discussion here.

In 2022, the worst year of the test period for equity markets, EGARCH is the best individual model with QLIKE of 0.2346, followed by GARCH(1,1) at 0.2579 and GJR-GARCH at 0.2597. The tree-based models perform notably worse: XGBoost at 0.3148 and LightGBM at 0.3429, with LightGBM ranking last among the seven. The ensemble lands in the middle at 0.2616. In the calmer 2023--2025 period, the ranking inverts: tree models lead (XGBoost 0.3693, LightGBM 0.3702), the ensemble is competitive at 0.3712, and GJR-GARCH (0.3740) leads the GARCH family but trails the trees, with HAR-RV last at 0.4568.

This inversion is the central nuanced finding of the paper. Tree-based ML models struggle in 2022 not because the 2022 vol level is unprecedented (the training data includes the 2008 financial crisis and the March 2020 COVID crash, both of which delivered higher RV than 2022) but because the immediate 2019--2021 validation period that drives early stopping and final feature importances was unusually calm. The trees' parameters are tuned to feature interactions that worked in the post-COVID-shock recovery, and 2022 is the first sustained departure from that calm regime. Parametric GARCH models, by contrast, have only three or four parameters fit on the full 2004--2021 history, which is dominated by long-run variance dynamics rather than recent regime tuning. In 2023--2025, when the RV level normalizes and intra-year regimes look more like the training distribution, the larger tree feature panel begins to pay off and tree models overtake.

The full-sample QLIKE is therefore a weighted average over two qualitatively different forecasting environments. A practitioner who reads only the full-sample ranking would conclude that GJR-GARCH is the best individual model, when in fact GJR-GARCH is third in 2022 and fourth in 2023--2025. The ensemble's robustness comes from precisely this regime-averaging property: it is rarely best in any one period, but never disastrously wrong either.
